% ============================================================================
%  CGM Rental — Informe Tecnico de la Pagina Web
%  Documentacion estructural completa para continuidad del proyecto
%  Desarrollado por: Omar Antonio Ccencho Barazorda
%  Cliente: CGM Rental
% ============================================================================
\documentclass[11pt,a4paper]{article}

% ── Paquetes ────────────────────────────────────────────────────────────────
\usepackage[utf8]{inputenc}
\usepackage[T1]{fontenc}
\usepackage[spanish,es-tabla]{babel}
\usepackage[a4paper,margin=2.4cm]{geometry}
\usepackage{xcolor}
\usepackage{graphicx}
\usepackage{titlesec}
\usepackage{enumitem}
\usepackage{tabularx}
\usepackage{longtable}
\usepackage{booktabs}
\usepackage{array}
\usepackage{listings}
\usepackage{fancyhdr}
\usepackage{tcolorbox}
\usepackage[hidelinks,colorlinks=true,linkcolor=cgmgreen,urlcolor=cgmgreen]{hyperref}

% ── Colores corporativos CGM ─────────────────────────────────────────────────
\definecolor{cgmgreen}{HTML}{004D3D}
\definecolor{cgmlime}{HTML}{C5E86C}
\definecolor{cgmgray}{HTML}{4A4A4A}
\definecolor{cgmgreenlight}{HTML}{E8F0EE}
\definecolor{codebg}{HTML}{F5F5F5}

% ── Logos ────────────────────────────────────────────────────────────────────
\graphicspath{{./}}
\def\cgmlogo{cropped-CGM_logo_P_Color_tagline_RGB-1-3hhh-270x270.png}
\def\cgmlogoheader{cgm-logo.png}

% ── Estilos de títulos ──────────────────────────────────────────────────────
\titleformat{\section}
  {\color{cgmgreen}\Large\bfseries}
  {\color{cgmlime}\rule[-0.4ex]{4pt}{1.4em}\hspace{6pt}\thesection}{0.6em}{}
\titleformat{\subsection}{\color{cgmgreen}\large\bfseries}{\thesubsection}{1em}{}
\titleformat{\subsubsection}{\color{cgmgray}\normalsize\bfseries}{\thesubsubsection}{1em}{}

% ── Encabezado y pie de página ──────────────────────────────────────────────
\pagestyle{fancy}
\fancyhf{}
\renewcommand{\headrulewidth}{0.4pt}
\renewcommand{\headrule}{\hbox to\headwidth{\color{cgmgreen}\leaders\hrule height \headrulewidth\hfill}}
\fancyhead[L]{%
  \raisebox{-0.35\height}{\includegraphics[height=22pt]{\cgmlogoheader}}}
\fancyhead[R]{\textcolor{cgmgray}{Informe Técnico}}
\fancyfoot[L]{}
\fancyfoot[C]{\textcolor{cgmgreen}{\textbf{\thepage}}}
\fancyfoot[R]{\textcolor{cgmgray}{\footnotesize cgmrental.com}}

% ── Listings para codigo ────────────────────────────────────────────────────
\lstdefinestyle{cgmcode}{
    backgroundcolor=\color{codebg},
    basicstyle=\ttfamily\footnotesize,
    keywordstyle=\color{cgmgreen}\bfseries,
    stringstyle=\color{cgmgray},
    commentstyle=\color{gray}\itshape,
    breaklines=true,
    showstringspaces=false,
    frame=single,
    rulecolor=\color{cgmlime},
    framesep=4pt,
    xleftmargin=8pt,
    aboveskip=10pt,
    belowskip=10pt,
}
\lstset{style=cgmcode}

% ── Cajas temáticas ─────────────────────────────────────────────────────────
\newtcolorbox{infobox}[1][]{
    colback=cgmlime!18,
    colframe=cgmgreen,
    boxrule=0.7pt, arc=3pt,
    left=10pt,right=10pt,top=8pt,bottom=8pt, #1
}
\newtcolorbox{warningbox}[1][]{
    colback=red!6,
    colframe=red!55,
    boxrule=0.7pt, arc=3pt,
    left=10pt,right=10pt,top=8pt,bottom=8pt, #1
}
\newtcolorbox{tipbox}[1][]{
    colback=cgmlime!12,
    colframe=cgmlime!80!cgmgreen,
    boxrule=0.7pt, arc=3pt,
    left=10pt,right=10pt,top=8pt,bottom=8pt, #1
}

% ============================================================================
\begin{document}

% ── Portada ──────────────────────────────────────────────────────────────────
\begin{titlepage}
\thispagestyle{empty}
\begin{center}
\vspace*{0.5cm}

\includegraphics[width=5cm]{\cgmlogo}\\[0.8cm]

\textcolor{cgmgreen}{\rule{\linewidth}{2.5pt}}\\[0.5cm]

{\color{cgmgreen}\Huge\bfseries Informe Técnico}\\[0.2cm]
{\color{cgmgreen}\Huge\bfseries Sitio Web Corporativo}\\[0.4cm]

\textcolor{cgmlime!70!black}{\rule{0.5\linewidth}{1pt}}\\[0.4cm]

{\color{cgmgray}\LARGE Documentación Estructural}\\[0.3cm]
{\color{cgmgreen}\Large CGM Rental --- Alquiler de Maquinaria Pesada}\\[0.6cm]

\textcolor{cgmgreen}{\rule{0.6\linewidth}{0.5pt}}\\[1.5cm]

\begin{tcolorbox}[
    colback=cgmgreenlight,
    colframe=cgmgreen,
    boxrule=0.8pt,
    arc=4pt,
    width=0.85\linewidth,
    left=14pt, right=14pt, top=12pt, bottom=12pt,
]
\centering
\textbf{\large\color{cgmgreen}Información del documento}\\[0.4cm]
\begin{tabular}{r l}
\textbf{Empresa:}      & CGM Rental \\
\textbf{Módulo:}       & Sitio Web Corporativo \\
\textbf{Repositorio:}  & \texttt{CGM\_RENTAL\_PAGE\_HOME} \\
\textbf{Fecha:}        & \today \\
\end{tabular}
\end{tcolorbox}

\vfill

{\color{cgmgray}\small Documento técnico --- Distribución interna}

\end{center}
\end{titlepage}

\begin{infobox}
\textbf{Proposito del documento.} Esta documentacion describe la arquitectura,
tecnologias, modelo de datos, rutas y funcionamiento general del sitio web
corporativo de CGM Rental. Esta pensada para que cualquier desarrollador o
equipo tecnico que herede el proyecto pueda comprender, mantener y extender
la plataforma sin depender del autor original.
\end{infobox}

\vspace{0.4cm}
\tableofcontents
\newpage

% ============================================================================
\section{Resumen Ejecutivo}
% ============================================================================

El sitio web de \textbf{CGM Rental} es una aplicacion web monolitica
desarrollada en Python con el framework \textbf{Flask}, que sirve como vitrina
corporativa, catalogo de productos y backoffice de administracion para las
operaciones de la empresa en \textbf{Peru} y \textbf{Argentina}.

\subsection{Objetivos del producto}
\begin{itemize}[leftmargin=*]
    \item Presentar la oferta de \textbf{alquiler y venta de maquinaria pesada}
        de CGM Rental en sus dos paises de operacion.
    \item Permitir al usuario navegar el catalogo por sector (Construccion,
        Mineria, Agricola, Energia) y por tipo de equipo.
    \item Generar \textbf{leads comerciales} a traves de Salesforce
        Web-to-Lead, con un formulario de contacto y un carrito-cotizador.
    \item Centralizar canales corporativos: \textbf{blog/novedades},
        \textbf{portal de proveedores}, \textbf{canal de denuncias},
        \textbf{certificaciones ISO}, seccion \textbf{Capacita} y
        \textbf{Leasing Operativo}.
    \item Adaptar la experiencia a cada pais (banderas, telefonos, redes,
        sucursales, productos visibles) de manera transparente.
    \item Proporcionar un \textbf{panel de administracion} para que el
        equipo de CGM Rental pueda gestionar productos, banners, blog,
        sucursales y datos de contacto sin intervencion de desarrollo.
\end{itemize}

\subsection{Cifras del proyecto}

\begin{center}
\begin{tabular}{l r}
\toprule
\textbf{Indicador} & \textbf{Valor} \\
\midrule
Lineas de codigo Python (\texttt{app.py})              & $\sim$1\,160 \\
Lineas de codigo Python (\texttt{admin/routes.py})     & $\sim$1\,370 \\
Plantillas Jinja2 publicas (\texttt{templates/pages/}) & 13 \\
Plantillas Jinja2 del admin (\texttt{templates/admin/}) & 11 \\
Productos en el catalogo (carpetas de imagenes)        & $\sim$125 \\
Paises soportados                                       & 2 (PE, AR) \\
Tablas SQLite principales                               & 5 \\
Tablas SQLite del panel admin                           & 4 \\
Lineas CSS personalizado (\texttt{cgm.css})            & $\sim$4\,000 \\
\bottomrule
\end{tabular}
\end{center}

% ============================================================================
\section{Stack Tecnologico}
% ============================================================================

\subsection{Backend}

\begin{longtable}{l l p{8cm}}
\toprule
\textbf{Componente} & \textbf{Version} & \textbf{Proposito} \\
\midrule
Python              & 3.x       & Lenguaje principal del backend. \\
Flask               & 3.1.0     & Framework web (rutas, sesiones, templating). \\
Werkzeug            & 3.1.3     & Servidor WSGI / utilidades HTTP. \\
flask-compress      & 1.15      & Compresion gzip/brotli automatica. \\
python-dotenv       & 1.0.1     & Carga de variables de entorno desde \texttt{.env}. \\
requests            & 2.32.3    & Cliente HTTP. \\
SQLite              & 3.x (WAL) & Base de datos en \texttt{data/cgm.db}. \\
Pillow              & ---       & Conversion automatica de imagenes a WebP. \\
msal                & ---       & Autenticacion Microsoft / Azure AD. \\
smtplib (stdlib)    & ---       & Envio de correos transaccionales. \\
\bottomrule
\end{longtable}

\subsection{Frontend}

\begin{longtable}{l l p{7.5cm}}
\toprule
\textbf{Recurso} & \textbf{Origen} & \textbf{Uso} \\
\midrule
Bootstrap           & 5.3 (local)  & Grilla, modales, navbar, formularios. \\
Bootstrap Icons     & 1.11 (local) & Iconografia general. \\
Swiper              & 11 (local)   & Carruseles (hero, productos relacionados). \\
Font Awesome Brands & 6.5.1 CDN    & Iconos de redes sociales. \\
Fuentes propias     & local        & \texttt{Estricta.otf}, \texttt{Estricta\_Black.otf}. \\
JavaScript          & vanilla      & \texttt{main.js}, \texttt{cart.js}, \texttt{cotizador.js}. \\
\bottomrule
\end{longtable}

\begin{infobox}
\textbf{Activos locales vs CDN.} Bootstrap, Bootstrap Icons y Swiper se
sirven desde \texttt{static/css/}, \texttt{static/js/} y \texttt{static/fonts/}
en lugar de CDN externos. Esto reduce dependencias externas, mejora la
estabilidad y permite trabajar en local sin conexion.
\end{infobox}

\subsection{Servicios integrados}

\begin{itemize}[leftmargin=*]
    \item \textbf{Salesforce Web-to-Lead.} Los formularios de cotizacion
        publicos envian datos directamente al CRM via
        \texttt{webto.salesforce.com/servlet/servlet.WebToLead}.
    \item \textbf{Azure Active Directory.} Autenticacion del panel admin
        usando MSAL.
    \item \textbf{Cloudflare.} CDN, proteccion DDoS y proveedor del header
        \texttt{CF-IPCountry} para deteccion automatica de pais.
    \item \textbf{Google Maps.} Embeds para sucursales.
    \item \textbf{YouTube.} Embeds de videos institucionales y entrevistas
        del blog.
\end{itemize}

\subsection{Infraestructura}

\begin{itemize}[leftmargin=*]
    \item \textbf{Servidor:} VPS en Azure con Nginx (reverse proxy) +
        Gunicorn (\texttt{cgmrental.service}).
    \item \textbf{CDN / WAF:} Cloudflare delante de Nginx.
    \item \textbf{Dominio:} \texttt{cgmrental.com}.
    \item \textbf{Despliegue:} manual via SSH + \texttt{git pull} +
        \texttt{systemctl restart cgmrental}.
\end{itemize}

% ============================================================================
\section{Estructura de Directorios}
% ============================================================================

\begin{lstlisting}[language=bash]
CGM_RENTAL_PAGE_HOME/
|-- app.py                    # Aplicacion Flask principal (~1160 lineas)
|-- database.py               # Conexion SQLite + schema + migraciones
|-- cache.py                  # Cache en memoria con TTL e invalidacion
|-- countries.py              # Defaults de Peru y Argentina (fallback)
|-- products.json             # Seed inicial del catalogo (~98 KB)
|-- requirements.txt          # Dependencias Python
|-- GEOIP_SETUP.md            # Documentacion de redireccion por pais
|-- .gitignore
|
|-- data/
|   `-- cgm.db                # Base de datos SQLite de produccion
|
|-- admin/                    # Panel de administracion (blueprint Flask)
|   |-- __init__.py
|   |-- auth.py               # Login con Azure AD (MSAL)
|   |-- products_sync.py      # Sincroniza BD <-> products.json
|   `-- routes.py             # Rutas y logica del panel (1370 lineas)
|
|-- templates/
|   |-- base.html             # Layout maestro
|   |-- _nav.html             # Nav superior + selector de pais (2 botones)
|   |-- _footer.html          # Footer completo (solo en /contacto/)
|   |-- _mini_footer.html     # Footer reducido para el resto del sitio
|   |
|   |-- pages/                # Plantillas publicas (13)
|   |   |-- home.html
|   |   |-- categoria.html
|   |   |-- producto.html
|   |   |-- ficha_tecnica.html
|   |   |-- carrito.html
|   |   |-- blog_list.html
|   |   |-- blog_post.html
|   |   |-- contacto.html
|   |   |-- gracias.html
|   |   |-- nosotros.html
|   |   |-- nuestros_locales.html
|   |   |-- leasing.html
|   |   |-- capacita.html
|   |   `-- canal_denuncias.html
|   |
|   `-- admin/                # Plantillas del panel admin (11)
|       |-- base.html
|       |-- login.html
|       |-- dashboard.html
|       |-- productos.html
|       |-- producto_form.html
|       |-- banners.html
|       |-- contacto.html
|       |-- sucursales.html
|       |-- blog_list.html
|       |-- blog_form.html
|       `-- historial.html
|
|-- static/
|   |-- css/
|   |   |-- cgm.css                  # Estilos principales (~4000 lineas)
|   |   |-- capacita.css             # Estilos de la pagina Capacita
|   |   |-- admin.css                # Estilos del panel admin
|   |   |-- bootstrap.min.css        # Bootstrap 5.3 local
|   |   |-- bootstrap-icons.min.css  # Bootstrap Icons local
|   |   `-- swiper-bundle.min.css    # Swiper 11 local
|   |
|   |-- js/
|   |   |-- main.js                  # Animaciones, nav, lazyload
|   |   |-- cart.js                  # Carrito y cotizador flotante
|   |   |-- cotizador.js             # Helpers del form de cotizacion
|   |   |-- bootstrap.bundle.min.js  # Bootstrap JS local
|   |   `-- swiper-bundle.min.js     # Swiper JS local
|   |
|   |-- fonts/
|   |   |-- Estricta.otf
|   |   |-- Estricta_Black.otf
|   |   |-- bootstrap-icons.woff
|   |   `-- bootstrap-icons.woff2
|   |
|   |-- products/             # 125 carpetas (una por producto):
|   |                         # 1.webp, 2.webp, 3.webp...
|   |
|   |-- images/
|   |   |-- banners/          # Banners organizados por subcarpetas:
|   |   |   |-- hero/         #   - Hero principal del home
|   |   |   |-- categorias/   #   - Cabeceras de pagina de categoria
|   |   |   |-- home-cards/   #   - Tarjetas del home (PE / AR)
|   |   |   `-- otros/        #   - Contacto, carrito, portal, etc.
|   |   |-- blog/             # Imagenes de posts
|   |   |-- logo/             # Logos en distintas resoluciones
|   |   |-- partners/         # Logos de clientes/partners (SVG)
|   |   |-- nosotros/         # Galerias de la seccion Nosotros
|   |   |-- fotos-capacita/   # Galeria de la seccion Capacita
|   |   |-- sig/              # Badges ISO
|   |   `-- contacto/         # Imagenes de la pagina de contacto
|   |
|   |-- docs/                 # PDFs: certificaciones ISO + brochures
|   |   `-- fichas/           # Fichas tecnicas de productos
|   |
|   `-- flags/                # SVGs de banderas para selector de pais
|
|-- docs/                     # Documentacion LaTeX del proyecto
|   |-- INFORME_TECNICO.tex
|   `-- MANUAL_PANEL_ADMIN.tex
|
|-- _no_usado/                # IGNORADO por git. Archivos viejos
|
`-- _archive/                 # IGNORADO por git. Backups historicos
\end{lstlisting}

% ============================================================================
\section{Modelo de Datos}
% ============================================================================

La base de datos es un fichero SQLite ubicado en \texttt{data/cgm.db}. Se abre
con tres pragmas optimizadas:

\begin{lstlisting}[language=Python]
conn.execute("PRAGMA journal_mode=WAL")    # mejor concurrencia
conn.execute("PRAGMA synchronous=NORMAL")  # mas rapido que FULL
conn.execute("PRAGMA foreign_keys=ON")
\end{lstlisting}

\begin{infobox}
\textbf{Pooling por request.} \texttt{database.py} mantiene una unica
conexion por request HTTP usando \texttt{flask.g}. Las llamadas dentro de
un request reciben un proxy (\texttt{\_ReqConn}) cuyo \texttt{close()} es
no-op. La conexion real se cierra en \texttt{teardown\_appcontext},
evitando el coste de abrir y cerrar SQLite varias veces por peticion.
\end{infobox}

\subsection{Migraciones}

No hay un sistema formal de migraciones (Alembic). \texttt{init\_db()} crea
las tablas si no existen y aplica \texttt{ALTER TABLE} envueltos en
\texttt{try/except} para agregar columnas a bases ya creadas.

\subsection{Tablas publicas (\texttt{init\_db})}

\subsubsection{\texttt{products}}
Catalogo principal de equipos.

\begin{longtable}{p{3.5cm} p{2cm} p{8cm}}
\toprule
\textbf{Columna} & \textbf{Tipo} & \textbf{Descripcion} \\
\midrule
\endhead
\texttt{id}                 & INTEGER PK  & Autoincremental. \\
\texttt{slug}               & TEXT UNIQUE & Identificador URL-friendly. \\
\texttt{nombre}             & TEXT        & Nombre comercial del equipo. \\
\texttt{marca}              & TEXT        & Marca del fabricante. \\
\texttt{descripcion}        & TEXT        & Features separadas por pipe (\texttt{|}). \\
\texttt{descripcion\_texto} & TEXT        & Texto largo descriptivo. \\
\texttt{ficha\_url}         & TEXT        & PDF local o URL externa. \\
\texttt{tags}               & TEXT        & \texttt{alquiler}, \texttt{venta} o ambas. \\
\texttt{tipo}               & TEXT        & Tipo de maquinaria. \\
\texttt{unidad}             & TEXT        & Sector. Acepta multivalor con pipe. \\
\texttt{imagen}             & TEXT        & Ruta principal (legacy). \\
\texttt{activo}             & INTEGER     & 1 visible, 0 oculto. \\
\texttt{show\_pe}           & INTEGER     & 1 visible en Peru. Default 1. \\
\texttt{show\_arg}          & INTEGER     & 1 visible en Argentina. Default 0. \\
\texttt{a\_solicitud}       & INTEGER     & 1 = equipo bajo solicitud. \\
\bottomrule
\end{longtable}

\subsubsection{\texttt{blog\_posts}}
Articulos y entrevistas del blog.

\begin{longtable}{p{3.2cm} p{2cm} p{8cm}}
\toprule
\textbf{Columna} & \textbf{Tipo} & \textbf{Descripcion} \\
\midrule
\endhead
\texttt{id}        & INTEGER PK  & Autoincremental. \\
\texttt{slug}      & TEXT UNIQUE & Identificador URL-friendly. \\
\texttt{titulo}    & TEXT        & Titulo del post. \\
\texttt{categoria} & TEXT        & \texttt{articulos}, \texttt{entrevistas}, etc. \\
\texttt{fecha}     & TEXT        & Formato \texttt{YYYY-MM-DD}. \\
\texttt{imagen}    & TEXT        & Ruta relativa en \texttt{static/images/}. \\
\texttt{extracto}  & TEXT        & Resumen para la card. \\
\texttt{contenido} & TEXT        & HTML completo del articulo. \\
\texttt{video\_url} & TEXT       & URL de YouTube opcional. \\
\texttt{activo}    & INTEGER     & 1 visible, 0 oculto. \\
\texttt{show\_pe}  & INTEGER     & 1 visible en Peru. Default 1. \\
\texttt{show\_arg} & INTEGER     & 1 visible en Argentina. Default 0. \\
\bottomrule
\end{longtable}

\subsubsection{\texttt{cotizaciones}}
Leads enviados al endpoint \texttt{/api/contacto}.

\begin{longtable}{p{3cm} p{2cm} p{9cm}}
\toprule
\textbf{Columna} & \textbf{Tipo} & \textbf{Descripcion} \\
\midrule
\endhead
\texttt{id}        & INTEGER PK & Autoincremental. \\
\texttt{nombre}    & TEXT       & Nombre del solicitante. \\
\texttt{empresa}   & TEXT       & Razon social. \\
\texttt{email}     & TEXT       & Correo. \\
\texttt{telefono}  & TEXT       & Telefono / WhatsApp. \\
\texttt{tipo}      & TEXT       & Tipo de consulta. \\
\texttt{mensaje}   & TEXT       & Mensaje libre. \\
\texttt{pais}      & TEXT       & \texttt{pe} o \texttt{ar}. \\
\texttt{productos} & TEXT       & Lista serializada del carrito. \\
\texttt{fecha}     & TEXT       & \texttt{datetime('now')}. \\
\bottomrule
\end{longtable}

\subsubsection{\texttt{proveedores}}
Registros del Portal de Proveedores.

\begin{longtable}{p{3cm} p{2cm} p{9cm}}
\toprule
\textbf{Columna} & \textbf{Tipo} & \textbf{Descripcion} \\
\midrule
\endhead
\texttt{id}, \texttt{empresa}, \texttt{ruc}, \texttt{contacto}, \texttt{email}, \texttt{telefono}, \texttt{rubro}, \texttt{descripcion} & INT/TEXT & Datos del proveedor. \\
\texttt{fecha} & TEXT & \texttt{datetime('now')}. \\
\bottomrule
\end{longtable}

\subsubsection{\texttt{denuncias}}
Reportes del Canal de Denuncias.

\begin{longtable}{p{3cm} p{2cm} p{9cm}}
\toprule
\textbf{Columna} & \textbf{Tipo} & \textbf{Descripcion} \\
\midrule
\endhead
\texttt{id}, \texttt{tipo}, \texttt{descripcion} & INT/TEXT & Categoria + descripcion. \\
\texttt{fecha} & TEXT & \texttt{datetime('now')}. \\
\bottomrule
\end{longtable}

\begin{infobox}
\textbf{Confidencialidad.} La tabla \texttt{denuncias} \textit{no almacena}
datos personales del denunciante. Estos datos viajan unicamente en el correo
HTML enviado al area de cumplimiento.
\end{infobox}

\subsection{Tablas del panel admin (\texttt{init\_admin\_tables})}

\subsubsection{\texttt{site\_config}}
Configuracion editable por pais.

\begin{longtable}{p{3.5cm} p{2cm} p{8cm}}
\toprule
\textbf{Columna} & \textbf{Tipo} & \textbf{Descripcion} \\
\midrule
\endhead
\texttt{id}            & INTEGER PK & Autoincremental. \\
\texttt{country\_code} & TEXT       & \texttt{pe} o \texttt{ar}. \\
\texttt{key}           & TEXT       & \texttt{telefono}, \texttt{whatsapp\_link}, \texttt{email}, \texttt{facebook}, \texttt{instagram}, \texttt{linkedin}, \texttt{youtube}, \texttt{direccion}, \texttt{ciudad\_principal}, \texttt{youtube\_video}. \\
\texttt{value}         & TEXT       & Valor del campo. \\
\multicolumn{3}{l}{\texttt{UNIQUE(country\_code, key)}} \\
\bottomrule
\end{longtable}

\subsubsection{\texttt{sucursales\_db}}
Sucursales gestionadas desde el admin.

\begin{longtable}{p{3.5cm} p{2cm} p{8cm}}
\toprule
\textbf{Columna} & \textbf{Tipo} & \textbf{Descripcion} \\
\midrule
\endhead
\texttt{id}            & INTEGER PK & Autoincremental. \\
\texttt{country\_code} & TEXT       & \texttt{pe} o \texttt{ar}. \\
\texttt{nombre}        & TEXT       & Nombre visible. \\
\texttt{tipo}          & TEXT       & \texttt{principal}, \texttt{sucursal}, \texttt{atencion}. \\
\texttt{direccion}     & TEXT       & Direccion fisica. \\
\texttt{lat, lng}      & REAL       & Coordenadas para Google Maps. \\
\texttt{maps\_url}     & TEXT       & Link directo en Google Maps. \\
\texttt{telefono}      & TEXT       & Opcional. \\
\texttt{orden}         & INTEGER    & Para ordenar manualmente. \\
\bottomrule
\end{longtable}

\subsubsection{\texttt{banners\_config}}
Configuracion de los banners del sitio.

\begin{longtable}{p{3.5cm} p{2cm} p{8cm}}
\toprule
\textbf{Columna} & \textbf{Tipo} & \textbf{Descripcion} \\
\midrule
\endhead
\texttt{slot}          & TEXT       & Identificador del slot. \\
\texttt{country\_code} & TEXT       & \texttt{*} (global), \texttt{pe} o \texttt{ar}. \\
\texttt{group\_name}   & TEXT       & Agrupacion visual (Hero, Home Cards, etc.). \\
\texttt{label}         & TEXT       & Nombre legible. \\
\texttt{description}   & TEXT       & Descripcion / contexto. \\
\texttt{filename}      & TEXT       & Ruta en \texttt{static/images/banners/}. \\
\texttt{orden}         & INTEGER    & Orden de presentacion. \\
\texttt{pages}         & TEXT       & JSON con lista de paginas. \\
\texttt{activo}        & INTEGER    & 0 / 1. \\
\bottomrule
\end{longtable}

\subsubsection{\texttt{audit\_log}}
Registro de toda accion administrativa.

\begin{longtable}{p{3cm} p{2cm} p{9cm}}
\toprule
\textbf{Columna} & \textbf{Tipo} & \textbf{Descripcion} \\
\midrule
\endhead
\texttt{id}      & INTEGER PK & Autoincremental. \\
\texttt{usuario} & TEXT       & Email del admin que ejecuto la accion. \\
\texttt{accion}  & TEXT       & Codigo de la accion. \\
\texttt{detalle} & TEXT       & Detalle libre. \\
\texttt{fecha}   & TEXT       & \texttt{datetime('now')}. \\
\bottomrule
\end{longtable}

% ============================================================================
\section{Sistema Multi-Pais}
% ============================================================================

El sitio sirve dos vitrinas separadas bajo el mismo dominio:

\begin{itemize}[leftmargin=*]
    \item \texttt{cgmrental.com/pe/} --- Peru (mercado principal).
    \item \texttt{cgmrental.com/ar/} --- Argentina.
\end{itemize}

\subsection{Dos niveles de configuracion}

\begin{enumerate}[leftmargin=*]
    \item \textbf{Defaults} en \texttt{countries.py}: diccionario hardcodeado
        con valores por defecto (banderas, flags de visibilidad, moneda) que
        se usa como fallback si la BD no tiene un valor.
    \item \textbf{Sobrescritura dinamica} desde la BD: las tablas
        \texttt{site\_config} y \texttt{sucursales\_db} permiten al panel
        admin modificar telefonos, redes, direccion principal y sucursales
        sin tocar codigo.
\end{enumerate}

\subsection{Selector de pais en la UI}

En el header del sitio se muestran \textbf{dos botones independientes}
(\texttt{Portal PE} / \texttt{Portal AR}) que enlazan a \texttt{/pe/} y
\texttt{/ar/} respectivamente. El boton del pais actual se resalta con borde
y fondo verde. Anteriormente esto era un dropdown; se cambio a botones
directos para reducir clicks y mejorar accesibilidad en movil.

\subsection{Deteccion automatica de pais}

La ruta raiz \texttt{/} ejecuta \texttt{\_detect\_country\_from\_ip()} que
lee el header \texttt{CF-IPCountry} de Cloudflare:

\begin{lstlisting}[language=Python]
def _detect_country_from_ip(_ip=None):
    cf = request.headers.get("CF-IPCountry", "").strip().upper()
    if cf and cf not in ("XX", "T1"):
        return {"PE": "pe", "AR": "ar"}.get(cf)
    return None
\end{lstlisting}

Cualquier otro pais cae a \texttt{DEFAULT\_COUNTRY = "pe"}.

\subsection{Filtrado por pais}

Los campos \texttt{show\_pe} y \texttt{show\_arg} en \texttt{products} y
\texttt{blog\_posts} controlan en que vitrina aparece cada elemento:

\begin{itemize}[leftmargin=*]
    \item Vitrina \texttt{/pe/}: muestra elementos con \texttt{show\_pe=1}.
    \item Vitrina \texttt{/ar/}: muestra elementos con \texttt{show\_arg=1}.
    \item Un mismo producto puede tener ambos flags si se vende en los dos
        paises (comun para mineria).
\end{itemize}

\subsection{Unidades multivalor}

El campo \texttt{unidad} acepta varios sectores separados por pipe
(\texttt{Construccion|Mediana Mineria}). El helper \texttt{\_unidad\_clause()}
en \texttt{app.py} construye el WHERE adecuado:

\begin{lstlisting}[language=Python]
def _unidad_clause(unidad, col="unidad"):
    return (
        f"({col} = ? OR {col} LIKE ? OR {col} LIKE ? OR {col} LIKE ?)",
        [unidad, f"{unidad}|%", f"%|{unidad}|%", f"%|{unidad}"]
    )
\end{lstlisting}

% ============================================================================
\section{Cache en Memoria}
% ============================================================================

El modulo \texttt{cache.py} implementa una capa de cache con TTL e
invalidacion explicita para datos que cambian con poca frecuencia:

\begin{itemize}[leftmargin=*]
    \item Configuracion del pais (\texttt{site\_config} + \texttt{sucursales\_db}).
    \item Banners del sitio (\texttt{banners\_config}).
\end{itemize}

\subsection{Estrategia de invalidacion}

\begin{itemize}[leftmargin=*]
    \item \textbf{Lectura:} cualquier request publico lee del cache.
    \item \textbf{Escritura desde admin:} cada vez que el panel admin
        guarda cambios, llama explicitamente a \texttt{invalidate()} para
        que el siguiente request publico lea fresco desde la BD.
\end{itemize}

Esto evita queries repetidas a SQLite en cada peticion y garantiza
simultaneamente que los cambios del admin se vean de inmediato.

% ============================================================================
\section{Mapa de Rutas Publicas}
% ============================================================================

\subsection{Paginas}

\begin{longtable}{p{6cm} p{3cm} p{5.5cm}}
\toprule
\textbf{Ruta} & \textbf{Metodo} & \textbf{Plantilla} \\
\midrule
\endhead
\texttt{/}                                            & GET & Redirect por GeoIP. \\
\texttt{/<country>/}                                  & GET & \texttt{pages/home.html} \\
\texttt{/<country>/nosotros/}                         & GET & \texttt{pages/nosotros.html} \\
\texttt{/<country>/nuestros-locales/}                 & GET & \texttt{pages/nuestros\_locales.html} \\
\texttt{/<country>/contacto/}                         & GET & \texttt{pages/contacto.html} \\
\texttt{/<country>/gracias/}                          & GET & \texttt{pages/gracias.html} \\
\texttt{/<country>/leasing-operativo/}                & GET & \texttt{pages/leasing.html} \\
\texttt{/<country>/capacita/}                         & GET & \texttt{pages/capacita.html} \\
\texttt{/<country>/canal-de-denuncias/}               & GET & \texttt{pages/canal\_denuncias.html} \\
\texttt{/<country>/portalproveedores/}                & GET & Redirect a portal externo. \\
\texttt{/<country>/novedades/}                        & GET & \texttt{pages/blog\_list.html} \\
\texttt{/<country>/novedades/page/<p>/}               & GET & Idem, paginado. \\
\texttt{/<country>/novedades/<slug>/}                 & GET & \texttt{pages/blog\_post.html} \\
\texttt{/<country>/categoria-producto/<path>/}        & GET & \texttt{pages/categoria.html} \\
\texttt{/<country>/producto/<slug>/}                  & GET & \texttt{pages/producto.html} \\
\texttt{/<country>/producto/<slug>/ficha/}            & GET & \texttt{pages/ficha\_tecnica.html} \\
\texttt{/<country>/carrito/}                          & GET & \texttt{pages/carrito.html} \\
\bottomrule
\end{longtable}

\subsection{SEO y utilitarias}

\begin{longtable}{p{5cm} p{2cm} p{8.5cm}}
\toprule
\textbf{Ruta} & \textbf{Metodo} & \textbf{Funcion} \\
\midrule
\endhead
\texttt{/robots.txt}    & GET & Politica para crawlers + apunta al sitemap. \\
\texttt{/sitemap.xml}   & GET & Sitemap dinamico por pais. \\
\bottomrule
\end{longtable}

\subsection{Endpoints API (JSON)}

\begin{longtable}{p{4.5cm} p{2cm} p{8.5cm}}
\toprule
\textbf{Ruta} & \textbf{Metodo} & \textbf{Funcion} \\
\midrule
\endhead
\texttt{/api/contacto}    & POST          & Guarda cotizacion + email. Valida 7 campos. \\
\texttt{/api/cart/add}    & POST          & Agrega un producto al carrito. \\
\texttt{/api/cart/qty}    & POST          & Actualiza cantidad de item. \\
\texttt{/api/cart}        & GET / DELETE  & Lee o elimina items del carrito. \\
\texttt{/api/denuncia}    & POST          & Persiste denuncia + email confidencial. \\
\texttt{/api/proveedor}   & POST          & Registra proveedor + email. \\
\bottomrule
\end{longtable}

\subsection{Diccionario \texttt{CATEGORIAS}}

En \texttt{app.py} se define un diccionario que mapea slugs URL a tuplas
\texttt{(tags, unidad, tipo, titulo)}. Es la fuente de verdad para los
filtros del catalogo. Cubre:

\begin{itemize}[leftmargin=*]
    \item 2 etiquetas principales: \texttt{alquiler}, \texttt{usados}.
    \item 4 unidades: Construccion, Mediana Mineria, Agricola, Energia.
    \item $\sim$25 tipos de maquinaria.
\end{itemize}

% ============================================================================
\section{Formularios y Validacion}
% ============================================================================

\subsection{Formularios de cotizacion (envio directo a Salesforce)}

Tanto \texttt{templates/pages/contacto.html} como
\texttt{templates/pages/carrito.html} envian sus datos directo al endpoint
de Salesforce Web-to-Lead via \texttt{POST}, no a la API interna. El payload
incluye campos con IDs de Salesforce (e.g.\ \texttt{00NUU000001AZFF},
\texttt{00NUU000001NBa5}).

\subsubsection{Campos obligatorios}

Todos los campos visibles de los dos formularios son \texttt{required}:

\begin{itemize}[leftmargin=*]
    \item Nombres y Apellidos, Razon Social, Email, Celular,
        RUC/DNI/CUIT, Pais, Departamento/Provincia.
    \item \textbf{Contacto:} ademas Equipo requerido.
    \item \textbf{Carrito:} ademas valida que el carrito tenga al menos un
        producto.
    \item Politica de privacidad: checkbox obligatorio.
\end{itemize}

\subsubsection{Validaciones especiales en JS}

\begin{itemize}[leftmargin=*]
    \item \textbf{Tipo de operacion} (checkboxes Alquiler / Compra): se
        valida via JS que al menos uno este marcado.
    \item \textbf{Otro Pais} (solo en contacto.html): se hace
        \texttt{required} dinamicamente solo cuando el select de Pais tiene
        valor \texttt{"Otro"}.
    \item \textbf{Carrito vacio}: bloqueo de submit si la tabla del carrito
        no tiene filas.
\end{itemize}

\subsection{Endpoint \texttt{/api/contacto}}

Aunque los formularios publicos van directo a Salesforce, el endpoint
interno se mantiene como respaldo. Valida 7 campos obligatorios
(\texttt{nombre}, \texttt{empresa}, \texttt{email}, \texttt{telefono},
\texttt{tipo}, \texttt{mensaje}, \texttt{pais}) y formato basico de email.
La respuesta de error incluye un campo \texttt{campos\_faltantes} con la
lista exacta de campos vacios.

\subsection{Canal de denuncias}

\texttt{templates/pages/canal\_denuncias.html} envia via AJAX a
\texttt{/api/denuncia}. Los datos del denunciante son \textbf{opcionales por
diseno} (anonimato). Solo son requeridos: nombre del denunciado, tipo de
denuncia y descripcion de los hechos.

% ============================================================================
\section{Carrito y Cotizador}
% ============================================================================

El carrito vive en la \textbf{sesion Flask} (cookie firmada con
\texttt{FLASK\_SECRET\_KEY}). No requiere registro de usuario.

\subsection{Separacion por pais}

Cada pais tiene su propio carrito independiente:

\begin{lstlisting}[language=Python]
def _cart_key(country=None):
    cc = (country or 'pe').lower()
    return f"cart_{cc}"   # cart_pe, cart_ar
\end{lstlisting}

\subsection{Cotizador flotante}

Definido en \texttt{base.html}, es un panel desplegable + boton que se
muestra en todas las paginas excepto en \texttt{/carrito/}. Permite al
usuario revisar los equipos acumulados sin abandonar la navegacion. Se
sincroniza con la sesion via \texttt{cart.js}.

\subsection{Estructura interna del carrito}

\begin{lstlisting}[language=Python]
{
    "slug":   "excavadora-john-deere-210g",
    "nombre": "Excavadora John Deere 210G",
    "imagen": "excavadora-john-deere-210g/1.webp",
    "tipo":   "Excavadora",
    "qty":    2
}
\end{lstlisting}

% ============================================================================
\section{Productos Relacionados con Prioridades}
% ============================================================================

En la ficha de cada producto (\texttt{/producto/<slug>/}) se muestran hasta
6 productos relacionados. El algoritmo aplica dos niveles de prioridad:

\begin{enumerate}[leftmargin=*]
    \item \textbf{Nivel 1.} Mismo \texttt{tipo} + misma \texttt{unidad} +
        mismo \texttt{tag}. Maximo 3.
    \item \textbf{Nivel 2.} Mismo \texttt{tag} + misma \texttt{unidad}, pero
        \texttt{tipo} distinto. Completa hasta 6.
\end{enumerate}

Ambos niveles aplican filtro \texttt{show\_pe/show\_arg} segun el pais y
usan \texttt{ORDER BY RANDOM()} para rotacion natural.

% ============================================================================
\section{Capacita: Centro de Entrenamiento}
% ============================================================================

La pagina \texttt{/capacita/} presenta los 4 programas de formacion de
operadores de maquinaria pesada. Los datos estan hardcodeados en la
plantilla.

\begin{longtable}{p{4.5cm} p{2.5cm} p{2.5cm} p{4cm}}
\toprule
\textbf{Programa} & \textbf{Duracion} & \textbf{Inversion} & \textbf{Modalidad} \\
\midrule
Formacion Basica         & 12 dias & S/. 3,250 & Presencial \\
Especializacion          & 10 dias & S/. 2,800 & Semi Presencial \\
Perfeccionamiento        & 9 dias  & S/. 2,800 & Semi Presencial \\
Actualizacion Virtual    & 6 dias  & S/. 500   & 100\% Virtual \\
\bottomrule
\end{longtable}

Todos los programas incluyen \textbf{Certificado de aprobacion a nombre de
CGM Rental} en su seccion de Beneficios.

% ============================================================================
\section{Emails Transaccionales}
% ============================================================================

La funcion \texttt{send\_email()} en \texttt{app.py} envia correos HTML via
SMTP plano (TLS en puerto 587). Las credenciales viven en variables de
entorno:

\begin{lstlisting}[language=bash]
SMTP_HOST=smtp.gmail.com
SMTP_PORT=587
SMTP_USER=...@cgmrental.com
SMTP_PASS=<app-password>
EMAIL_DESTINO=contacto@cgmrental.com
\end{lstlisting}

Las plantillas HTML estan \textbf{embebidas en el codigo Python}. Cada
formulario tiene su propia plantilla: cotizacion, denuncia (estilo corporativo
verde / lima), proveedor.

\begin{warningbox}
Si \texttt{SMTP\_USER} o \texttt{SMTP\_PASS} no estan configurados,
\texttt{send\_email()} devuelve \texttt{False} silenciosamente. La cotizacion
si se guarda en la base de datos, pero no llega al correo de ventas.
Verificar el \texttt{.env} en cada despliegue.
\end{warningbox}

% ============================================================================
\section{Frontend}
% ============================================================================

\subsection{Plantilla maestra (\texttt{base.html})}

Define:
\begin{itemize}[leftmargin=*]
    \item \texttt{<head>}: metatags, favicons, preload de fuentes, CSS
        locales (Bootstrap, Bootstrap Icons, Swiper) + CSS custom
        (\texttt{cgm.css}, \texttt{capacita.css}) + FontAwesome brands
        desde CDN.
    \item Bloques Jinja: \texttt{title}, \texttt{meta\_desc},
        \texttt{og\_title}, \texttt{og\_desc}, \texttt{extra\_head},
        \texttt{content}, \texttt{extra\_scripts}.
    \item Inclusion del nav (\texttt{\_nav.html}) y del footer (
        \texttt{\_footer.html} si la URL contiene \texttt{/contacto/}, sino
        \texttt{\_mini\_footer.html}).
    \item Cotizador flotante con badge de cantidad.
    \item Boton de WhatsApp flotante.
    \item Variable global: \texttt{window.CGM\_COUNTRY = "<pe|ar>"}.
\end{itemize}

\subsection{Context processor}

\texttt{app.py} inyecta variables globales en todas las plantillas:

\begin{lstlisting}[language=Python]
@app.context_processor
def inject_globals():
    path_parts = request.path.strip('/').split('/')
    cc = path_parts[0] if path_parts and path_parts[0] in COUNTRIES else 'pe'
    return {
        "cart_count": cart_count(cc),
        "COUNTRIES":  COUNTRIES,
        "PARTNERS":   PARTNERS,
    }
\end{lstlisting}

\subsection{JavaScript}

Tres archivos en vanilla JS, sin frameworks:

\begin{itemize}[leftmargin=*]
    \item \texttt{main.js}: Swipers, comportamiento del nav, animaciones,
        lazyload, menu mobile.
    \item \texttt{cart.js}: llamadas a \texttt{/api/cart/*}, actualiza badge,
        sincroniza panel del cotizador.
    \item \texttt{cotizador.js}: helper para el envio del formulario.
\end{itemize}

% ============================================================================
\section{Panel de Administracion}
% ============================================================================

Este informe describe el panel desde un enfoque tecnico. Para la guia de
uso operativo, ver el documento separado
\textit{Manual del Panel de Administracion}.

\subsection{Arquitectura}

\begin{itemize}[leftmargin=*]
    \item \textbf{Blueprint Flask:} montado en \texttt{/admin}, registrado
        en \texttt{app.py}:
\begin{lstlisting}[language=Python]
from admin.routes import admin_bp
app.register_blueprint(admin_bp)
\end{lstlisting}
    \item \textbf{Modulos:}
    \begin{itemize}
        \item \texttt{admin/routes.py} (1370 lineas) --- todas las rutas y
            logica del panel.
        \item \texttt{admin/auth.py} --- autenticacion Azure AD con MSAL.
        \item \texttt{admin/products\_sync.py} --- sincronizacion bidireccional
            entre la BD y \texttt{products.json}.
        \item \texttt{admin/\_\_init\_\_.py}.
    \end{itemize}
    \item \textbf{Plantillas:} 11 templates en \texttt{templates/admin/}.
    \item \textbf{CSS:} \texttt{static/css/admin.css}.
    \item \textbf{Tablas:} \texttt{audit\_log}, \texttt{site\_config},
        \texttt{sucursales\_db}, \texttt{banners\_config}.
\end{itemize}

\subsection{Rutas principales}

\begin{longtable}{p{6cm} p{2cm} p{6cm}}
\toprule
\textbf{Ruta} & \textbf{Metodo} & \textbf{Funcion} \\
\midrule
\endhead
\texttt{/admin/login}                       & GET           & Pantalla de login Azure AD. \\
\texttt{/admin/auth/callback}               & GET           & Callback OAuth2. \\
\texttt{/admin/logout}                      & GET           & Cierre de sesion. \\
\texttt{/admin/dashboard}                   & GET           & KPIs y graficos. \\
\texttt{/admin/dashboard/data}              & GET           & Datos AJAX filtrados por pais. \\
\texttt{/admin/productos}                   & GET           & Listado paginado de productos. \\
\texttt{/admin/productos/nuevo}             & GET/POST      & Crear producto. \\
\texttt{/admin/productos/<id>/editar}       & GET/POST      & Editar producto. \\
\texttt{/admin/productos/<id>/toggle}       & POST          & Activar / desactivar. \\
\texttt{/admin/productos/sync-json}         & POST          & Reconstruir \texttt{products.json}. \\
\texttt{/admin/upload}                      & POST          & Subir imagen (convierte a WebP). \\
\texttt{/admin/banners}                     & GET           & Listar banners + huerfanos. \\
\texttt{/admin/banners/replace}             & POST          & Reemplazar imagen de slot. \\
\texttt{/admin/banners/country-upload}      & POST          & Crear version por pais. \\
\texttt{/admin/banners/country-remove}      & POST          & Eliminar version por pais. \\
\texttt{/admin/contacto}                    & GET/POST      & Editar \texttt{site\_config}. \\
\texttt{/admin/sucursales}                  & GET           & Listado de sucursales. \\
\texttt{/admin/sucursales/guardar}          & POST          & Crear / editar sucursal. \\
\texttt{/admin/blog}                        & GET           & Listado de posts. \\
\texttt{/admin/blog/nuevo}                  & GET/POST      & Crear post. \\
\texttt{/admin/blog/<id>/editar}            & GET/POST      & Editar post. \\
\texttt{/admin/historial}                   & GET           & Audit log paginado. \\
\bottomrule
\end{longtable}

\subsection{Autenticacion}

\begin{itemize}[leftmargin=*]
    \item OAuth2 contra Azure AD usando \textbf{MSAL} (Microsoft
        Authentication Library).
    \item Variables de entorno: \texttt{AZURE\_CLIENT\_ID},
        \texttt{AZURE\_CLIENT\_SECRET}, \texttt{AZURE\_TENANT\_ID},
        \texttt{AZURE\_REDIRECT\_URI}.
    \item Whitelist de correos en \texttt{ADMIN\_ALLOWED\_EMAILS}
        (lista separada por comas). Solo los emails de esta lista pueden
        entrar al panel.
    \item Sesion en cookie Flask firmada.
\end{itemize}

\subsection{Procesamiento de imagenes}

Todas las imagenes subidas pasan por Pillow:

\begin{itemize}[leftmargin=*]
    \item Si la extension es \texttt{webp}: se guarda tal cual.
    \item Si es \texttt{jpg/jpeg/png}: se convierte a WebP con calidad 88,
        respetando transparencia para RGBA/LA/P.
    \item Imagenes de producto se renumeran automaticamente
        (\texttt{1.webp}, \texttt{2.webp}, ...).
    \item Tamano maximo: 10 MB.
    \item Banners conservan su \texttt{filename} al reemplazarse.
\end{itemize}

% ============================================================================
\section{SEO}
% ============================================================================

\subsection{\texttt{robots.txt}}
Servido dinamicamente en \texttt{/robots.txt}. Permite todo y apunta al
sitemap.

\subsection{\texttt{sitemap.xml}}
Generado en cada request en \texttt{/sitemap.xml}. Incluye:

\begin{itemize}[leftmargin=*]
    \item Paginas estaticas por pais (home, nosotros, contacto, leasing,
        novedades, capacita, carrito, canal de denuncias).
    \item Todas las categorias de productos activas por pais.
    \item Todos los productos activos por pais (filtrados con
        \texttt{show\_pe}/\texttt{show\_arg}).
\end{itemize}

% ============================================================================
\section{Variables de Entorno}
% ============================================================================

\begin{longtable}{p{5cm} p{2cm} p{8cm}}
\toprule
\textbf{Variable} & \textbf{Default} & \textbf{Uso} \\
\midrule
\endhead
\texttt{FLASK\_SECRET\_KEY}         & \texttt{cgm-dev-secret} & Firma de cookies de sesion. \textbf{Obligatoria en produccion}. \\
\texttt{FLASK\_DEBUG}               & \texttt{True}           & Modo debug. \texttt{False} en produccion. \\
\texttt{SMTP\_HOST}                 & \texttt{smtp.gmail.com} & Servidor de correo. \\
\texttt{SMTP\_PORT}                 & \texttt{587}            & Puerto SMTP. \\
\texttt{SMTP\_USER}                 & ---                     & Cuenta de envio. \\
\texttt{SMTP\_PASS}                 & ---                     & Password / app-password. \\
\texttt{EMAIL\_DESTINO}             & \texttt{contacto@...}   & Destinatario de cotizaciones. \\
\texttt{AZURE\_CLIENT\_ID}          & ---                     & App ID de Azure AD. \\
\texttt{AZURE\_CLIENT\_SECRET}      & ---                     & Secret de Azure AD. \\
\texttt{AZURE\_TENANT\_ID}          & ---                     & Tenant ID. \\
\texttt{AZURE\_REDIRECT\_URI}       & \texttt{.../callback}   & Callback OAuth2. \\
\texttt{ADMIN\_ALLOWED\_EMAILS}     & ---                     & Whitelist (comas). \\
\bottomrule
\end{longtable}

% ============================================================================
\section{Despliegue}
% ============================================================================

\subsection{Local (desarrollo)}

\begin{lstlisting}[language=bash]
git clone <repo>
cd CGM_RENTAL_PAGE_HOME
python -m venv venv
source venv/bin/activate     # Linux/Mac
# .\venv\Scripts\Activate.ps1  # Windows PowerShell
pip install -r requirements.txt
cp .env.example .env          # editar credenciales SMTP / Azure
python app.py                 # http://localhost:5000
\end{lstlisting}

\subsection{Produccion (Azure VPS)}

\begin{lstlisting}[language=bash]
cd /var/www/cgmrental
git pull origin main
source venv/bin/activate
pip install -r requirements.txt
sudo systemctl restart cgmrental
\end{lstlisting}

El servicio \texttt{cgmrental.service} corre Gunicorn:

\begin{lstlisting}[language=bash]
gunicorn -w 4 -b 0.0.0.0:8015 app:app --daemon
\end{lstlisting}

Nginx hace reverse proxy desde \texttt{:443} hacia \texttt{:8015} y propaga
\texttt{X-Forwarded-For} y \texttt{CF-IPCountry} de Cloudflare.

\subsection{Carpetas ignoradas por Git}

El \texttt{.gitignore} excluye del repositorio:

\begin{itemize}[leftmargin=*]
    \item \texttt{.env}, \texttt{.env.local} --- secretos.
    \item \texttt{.claude/}, \texttt{.vscode/}, \texttt{.idea/} --- editor.
    \item \texttt{\_\_pycache\_\_/}, \texttt{*.pyc} --- bytecode Python.
    \item \texttt{\_archive/}, \texttt{\_no\_usado/} --- archivos historicos
        y descartados.
    \item \texttt{*.db}, \texttt{*.sqlite} --- BDs locales, \textbf{excepto}
        \texttt{data/cgm.db} (incluida con \texttt{!data/cgm.db}).
    \item \texttt{*.xlsx}, \texttt{*.sql} --- artefactos de importacion.
\end{itemize}

% ============================================================================
\section{Historial Tecnico Relevante}
% ============================================================================

\subsection{Migracion a SQL Server (revertida)}
Hubo un intento de migrar la base de datos de SQLite a Azure SQL Server. La
migracion fue revertida porque \textbf{SQLite resultaba mas rapido} para los
patrones de lectura del sitio. La base de datos sigue siendo SQLite local
en \texttt{data/cgm.db}.

\subsection{Selector de pais: dropdown $\rightarrow$ dos botones}
El header anteriormente tenia un boton \texttt{Portal PE} con dropdown que
listaba paises. Se reemplazo por dos botones independientes lado a lado
(\texttt{Portal PE} y \texttt{Portal AR}) para reducir clicks y mejorar
accesibilidad en movil.

\subsection{Reorganizacion de banners}
Los banners pasaron de estar todos en \texttt{static/images/banners/} a una
estructura por subcarpetas: \texttt{hero/}, \texttt{categorias/},
\texttt{home-cards/}, \texttt{otros/}, cada una con sub-carpetas
\texttt{pe/} y \texttt{ar/} para versiones especificas por pais.

\subsection{Migracion CDN $\rightarrow$ locales}
Bootstrap, Bootstrap Icons y Swiper se migraron de CDN a archivos locales
en \texttt{static/} para mayor estabilidad y autonomia.

\subsection{Endurecimiento de validacion de formularios}
Los formularios de cotizacion (\texttt{contacto.html}, \texttt{carrito.html})
ahora exigen \textbf{todos los campos visibles} como obligatorios, validan
que al menos uno de Alquiler/Compra este marcado, y que el carrito no este
vacio. El endpoint \texttt{/api/contacto} valida 7 campos en lugar de los
3 originales.

% ============================================================================
\section{Consideraciones de Seguridad}
% ============================================================================

\begin{itemize}[leftmargin=*]
    \item \textbf{SQL Injection:} todas las queries usan parametros
        \texttt{?}. Las concatenaciones con f-strings solo emplean valores
        de una whitelist.
    \item \textbf{CSRF:} los formularios publicos van directo a Salesforce
        (origin externo), por lo que no llevan CSRF token propio. Las rutas
        POST internas consumen JSON sin estado autenticado, por lo que el
        riesgo CSRF es bajo.
    \item \textbf{XSS:} Jinja2 escapa por defecto. El campo
        \texttt{blog\_posts.contenido} se renderiza con \texttt{|safe} ---
        debe venir de admins de confianza unicamente.
    \item \textbf{Cookie de sesion:} firmada con \texttt{FLASK\_SECRET\_KEY}.
        En produccion debe ser un valor unico y aleatorio.
    \item \textbf{Panel admin:} doble capa de proteccion --- Azure AD +
        whitelist de emails.
    \item \textbf{Subida de archivos:} solo desde el panel admin, con
        validacion de extensiones y tamano $\leq$ 10 MB.
\end{itemize}

% ============================================================================
\section{Hoja de Ruta Sugerida}
% ============================================================================

\begin{enumerate}[leftmargin=*]
    \item Migracion formal de schema con Alembic en lugar de
        \texttt{ALTER TABLE} con \texttt{try/except}.
    \item Vista admin para gestionar cotizaciones / proveedores / denuncias
        recibidas (actualmente solo se guardan en BD y se envian por email).
    \item Cola asincronica (\texttt{Celery} + \texttt{Redis}) para email,
        para no bloquear el request del usuario.
    \item Sistema de roles dentro del admin (actualmente cualquier email
        autorizado tiene acceso total).
    \item Tests automatizados (\texttt{pytest}).
    \item Logging estructurado en archivos rotativos.
    \item Migracion del HTML inline de emails a templates Jinja separados.
\end{enumerate}

% ============================================================================
\section{Glosario}
% ============================================================================

\begin{description}[leftmargin=*]
    \item[\texttt{slug}] Identificador URL-friendly de un producto o post.
        Ejemplo: \texttt{excavadora-john-deere-210g}.
    \item[\texttt{tag}] Categoria comercial: \texttt{alquiler} o
        \texttt{venta} (alias \texttt{usados}).
    \item[\texttt{unidad}] Sector economico: Construccion, Mediana Mineria,
        Agricola, Energia.
    \item[\texttt{tipo}] Tipo de equipo dentro de un sector.
    \item[\texttt{show\_pe / show\_arg}] Flags binarios que indican en que
        vitrina aparece el producto/post.
    \item[\texttt{slot}] Posicion fija de un banner en el diseno.
    \item[\texttt{huerfano}] Archivo en \texttt{static/images/banners/} que
        no esta registrado en \texttt{banners\_config}.
    \item[\texttt{ficha\_url}] Ruta a la ficha tecnica PDF.
\end{description}

% ============================================================================
\section{Despliegue en Producción}
% ============================================================================

\subsection{Ubicación en el servidor}

El proyecto se encuentra en el servidor de producción bajo la ruta:

\begin{lstlisting}
/var/www/cgmrental/
\end{lstlisting}

El usuario del sistema operativo que ejecuta la aplicación es
\texttt{admincgm} en el host \texttt{mq-vtl-prd-cgm}.

\subsection{Servidor de aplicación}

La aplicación corre con \textbf{Gunicorn} (Green Unicorn), un servidor
WSGI de producción para Python. Configuración activa:

\begin{itemize}[leftmargin=*]
    \item \textbf{Workers:} 4 procesos paralelos (\texttt{-w 4}).
    \item \textbf{Bind:} \texttt{0.0.0.0:8015} (puerto 8015, todas las
        interfaces).
    \item \textbf{App:} \texttt{app:app} (módulo \texttt{app.py},
        instancia \texttt{app}).
\end{itemize}

\subsection{Script de arranque}

El archivo \texttt{start.sh} crea una sesión \textbf{screen} llamada
\texttt{cgmrental} en segundo plano y dentro de ella lanza Gunicorn:

\begin{lstlisting}[language=bash]
#!/bin/bash
screen -dmS cgmrental bash -c '
cd /var/www/cgmrental
source venv/bin/activate
gunicorn -w 4 -b 0.0.0.0:8015 app:app
exec bash
'
\end{lstlisting}

El uso de \texttt{screen} permite que el proceso persista aunque se
cierre la sesión SSH.

\subsection{Servicio systemd}

La aplicación se registra como servicio del sistema operativo en
\texttt{/etc/systemd/system/cgmrental.service}, lo que garantiza que
arranque automáticamente al reiniciar el servidor:

\begin{lstlisting}
[Unit]
Description=CGM Rental Web
After=network.target

[Service]
User=admincgm
Type=forking
RemainAfterExit=yes
ExecStart=/usr/bin/screen -dmS cgmrental /var/www/cgmrental/start.sh
ExecStop=/usr/bin/screen -S cgmrental -X quit

[Install]
WantedBy=multi-user.target
\end{lstlisting}

\subsection{Comandos de administración del servicio}

\begin{longtable}{p{7cm} p{7cm}}
\toprule
\textbf{Acción} & \textbf{Comando} \\
\midrule
\endhead
Iniciar el servicio      & \texttt{sudo systemctl start cgmrental} \\
Detener el servicio      & \texttt{sudo systemctl stop cgmrental} \\
Reiniciar el servicio    & \texttt{sudo systemctl restart cgmrental} \\
Ver estado               & \texttt{sudo systemctl status cgmrental} \\
Ver sesión screen activa & \texttt{screen -ls} \\
Conectarse a la consola  & \texttt{screen -r cgmrental} \\
Salir sin cerrar screen  & \texttt{Ctrl+A} luego \texttt{D} \\
\bottomrule
\end{longtable}

\begin{tipbox}
\textbf{Después de subir cambios al servidor}, siempre ejecutar
\texttt{sudo systemctl restart cgmrental} para que Flask cargue el
código actualizado.
\end{tipbox}

\vspace{1cm}
\begin{tcolorbox}[
    colback=cgmgreen,
    colframe=cgmlime,
    boxrule=1pt,
    arc=4pt,
    coltext=white,
    left=14pt, right=14pt, top=12pt, bottom=12pt,
]
\begin{center}
\textbf{\Large Documento de uso interno}\\[0.3cm]
Este informe es de distribución restringida al equipo técnico de CGM Rental.\\[0.2cm]
Para consultas sobre la arquitectura o el servidor, comunicarse con\\
el área de sistemas o tecnología de la empresa.
\end{center}
\end{tcolorbox}

\end{document}
